\section{Top-level AXI4 interface}
\label{sec:axi4top}

CIF presents a standard AXI4 slave interface to its client(s).  The table
below cross-references the channel signals against the widths the current
architecture documents establish (\texttt{ADDR\_W=32}, \texttt{AXID\_W=6},
\texttt{LEN\_W=8}, from the Request Validator's Design Parameters).  Widths
marked AXI4-fixed follow the AXI4 protocol specification itself and are not
CIF-specific parameters.  The write and read data paths are \texttt{512}
bits: one 64-B AXI data beat per MC request packet (\S\ref{sec:identity},
\S\ref{sec:respformat}).

\begin{longtable}{p{1.3cm}p{3.2cm}p{1.6cm}p{9.5cm}}
\toprule
\textbf{Channel} & \textbf{Signal} & \textbf{Width} & \textbf{Notes}\\
\midrule
AW & \texttt{AWVALID/AWREADY} & 1/1 & Write address handshake\\
   & \texttt{AWADDR}          & 32  & \texttt{ADDR\_W}\\
   & \texttt{AWID}            & 6   & \texttt{AXID\_W}\\
   & \texttt{AWLEN}           & 8   & \texttt{LEN\_W}; raw AXI burst length. Given the 64-B-beat contract, \texttt{exp\_pkts = AWLEN + 1} (\S\ref{sec:identity})\\
   & \texttt{AWSIZE/AWBURST}  & 3/2 & AXI4-fixed; also feed the write 64\,B line footprint calc (\S\ref{sec:samelinehaz})\\
\midrule
W  & \texttt{WVALID/WREADY}   & 1/1 & Write data handshake\\
   & \texttt{WDATA}           & 512 & One 64-B AXI data beat per MC request packet (\S\ref{sec:identity})\\
   & \texttt{WSTRB}           & 64  & AXI4-fixed (one bit per \texttt{WDATA} byte); carried verbatim as the \texttt{wmask} field on \texttt{wdata\_afifo} (\S\ref{sec:cifmc})\\
   & \texttt{WLAST}           & 1   & AXI4-fixed\\
\midrule
B  & \texttt{BVALID/BREADY}   & 1/1 & Write response handshake\\
   & \texttt{BID}             & 6   & \texttt{AXID\_W}; from the retiring transaction's Slot FIFO context, not stored in the Metadata RAM entry\\
   & \texttt{BRESP}           & 2   & AXI4-fixed; driven from the write ROB's aggregate \texttt{resp\_status}\\
\midrule
AR & \texttt{ARVALID/ARREADY} & 1/1 & Read address handshake\\
   & \texttt{ARADDR}          & 32  & \texttt{ADDR\_W}\\
   & \texttt{ARID}            & 6   & \texttt{AXID\_W}\\
   & \texttt{ARLEN}           & 8   & \texttt{LEN\_W}; raw AXI burst length. Given the 64-B-beat contract, \texttt{exp\_pkts = ARLEN + 1} (\S\ref{sec:identity})\\
   & \texttt{ARSIZE/ARBURST}  & 3/2 & AXI4-fixed; also feed the read footprint used by the same-line hazard check (\S\ref{sec:samelinehaz})\\
\midrule
R  & \texttt{RVALID/RREADY}   & 1/1 & Read data handshake\\
   & \texttt{RDATA}           & 512 & One 64-B AXI data beat per MC request packet (\S\ref{sec:identity})\\
   & \texttt{RID}             & 6   & \texttt{AXID\_W}; from the retiring transaction's Slot FIFO context\\
   & \texttt{RRESP}           & 2   & AXI4-fixed; per-beat value from the read ROB's serialized status\\
   & \texttt{RLAST}           & 1   & Asserted only on the final AXI beat of the transaction (\texttt{PACKET\_NUMBER == exp\_pkts - 1}), not on every packet\\
\bottomrule
\caption{Top-level AXI4 slave interface, cross-referenced to CIF-internal handling}
\end{longtable}

\section{Interface Signal Table}
\label{sec:interfaces}

This table records the consolidated interfaces after the read/write ROB
refactor.  Names describe the frozen connectivity where prior documents did
not fix a signal name.

\subsection{AXI request acceptance to the selected ROB}
\begin{longtable}{p{3.3cm}p{1cm}p{3.4cm}p{7cm}}
\toprule
\textbf{Signal} & \textbf{Width} & \textbf{Path} & \textbf{Description}\\
\midrule
\texttt{alloc\_req} & 1 & Accept logic $\rightarrow$ RD/WR ROB & Allocate one transaction slot on accepted AR/AW request.\\
\texttt{AXID} & 6 & Accept logic $\rightarrow$ RD/WR ROB & Selects the per-AXID Slot FIFO.\\
\texttt{ROB\_INDEX} & 8 & RD/WR ROB $\rightarrow$ Splitter & Free-List result retained for the transaction.\\
\texttt{exp\_pkts} & 7 & Splitter $\rightarrow$ RD/WR ROB & Number of 64-B MC packets for this transaction, 1--64; initializes the ROB's \texttt{exp\_pkts}.\\
\texttt{wr\_line\_base} & 26 & Accept logic $\rightarrow$ WR ROB & Start 64\,B line of the write footprint; written at allocation (\S\ref{sec:samelinehaz}).\\
\texttt{wr\_line\_count} & 7 & Accept logic $\rightarrow$ WR ROB & Contiguous 64\,B line span of the write, 1--64 (\S\ref{sec:samelinehaz}).\\
\texttt{HAZ\_HOLD} & 1 & Same-line hazard gate $\rightarrow$ Request Validator & Read path: head read overlaps an occupied \texttt{ROB\_WR} footprint; withholds \texttt{alloc\_req} (\S\ref{sec:samelinehaz}).\\
\texttt{free\_req} & --- & WR ROB $\rightarrow$ same-line hazard gate & Existing write-retirement pulse; also re-arms the held-read recheck (\S\ref{sec:samelinehaz}).\\
\bottomrule
\end{longtable}

\subsection{Burst Splitter to Packet Formation}
\begin{longtable}{p{3.3cm}p{1cm}p{3.4cm}p{7cm}}
\toprule
\textbf{Signal} & \textbf{Width} & \textbf{Path} & \textbf{Description}\\
\midrule
\texttt{emit\_valid/emit\_ready} & 1/1 & Splitter $\leftrightarrow$ Packet Formation & Handshake for the current 64-B MC packet.\\
\texttt{emit\_addr} & 32 & Splitter $\rightarrow$ Packet Formation & 64-B-granular start address of this packet's data beat.\\
\texttt{PACKET\_NUMBER} & 6 & Splitter $\rightarrow$ Packet Formation & Current \texttt{pkt\_idx} (\texttt{0..exp\_pkts-1}), assigned at packet emission.\\
\texttt{GLOBAL\_TAG} & 14 & Packet Formation $\rightarrow$ \texttt{req\_afifo} & \texttt{\{ROB\_INDEX[7:0],PACKET\_NUMBER[5:0]\}}.\\
\bottomrule
\end{longtable}

\noindent Packet Formation applies the CIF-side address map/hash to
\texttt{emit\_addr}, producing the packet's \texttt{daddr} field
(\S\ref{sec:cifmc}); the raw byte address is not carried to MC.

\subsection{MC Core completion interface}

Response transport is the valid-only \texttt{resp\_afifo} (\S\ref{sec:cifmc});
CIF always accepts, MC self-throttles via a per-read reserved slot.  MC Core
supplies exactly one response per accepted 64-B request packet
(\S\ref{sec:cifmc-req}).

\begin{longtable}{p{3.3cm}p{1cm}p{3.4cm}p{7cm}}
\toprule
\textbf{Signal} & \textbf{Width} & \textbf{Path} & \textbf{Description}\\
\midrule
\texttt{resp\_type} & 1 & MC Core $\rightarrow$ RD/WR ROB & \texttt{0 = RD\_DATA} (route to read ROB), \texttt{1 = WR\_ACK} (route to write ROB).\\
\texttt{GLOBAL\_TAG} & 14 & MC Core $\rightarrow$ RD/WR ROB & Echoed request tag; \texttt{[13:6]} = \texttt{ROB\_INDEX}, \texttt{[5:0]} = \texttt{PACKET\_NUMBER}.\\
\texttt{status} & 4 & MC Core $\rightarrow$ RD/WR ROB & \texttt{0 = OK}, \texttt{1 = SLVERR}, \texttt{2 = DECERR}; codes \texttt{3..15} reserved $\to$ \texttt{DECERR} (\S\ref{sec:cifmc}).\\
\texttt{rdata} & 512 & MC Core $\rightarrow$ Read ROB & Read data payload, one 64-B unit; valid only when \texttt{resp\_type = RD\_DATA} (\S\ref{sec:respformat}).\\
\bottomrule
\end{longtable}

\subsection{Retirement interfaces}
\begin{longtable}{p{3.3cm}p{1cm}p{3.4cm}p{7cm}}
\toprule
\textbf{Signal} & \textbf{Width} & \textbf{Path} & \textbf{Description}\\
\midrule
\texttt{merge\_ready} & 1 & Merge Logic $\rightarrow$ Read ROB & Accepts an emitted AXI read beat (one per 64-B MC packet).\\
\texttt{rob\_last} & 1 & Read ROB $\rightarrow$ Merge Logic & Asserted with the final serialized read beat (\texttt{PACKET\_NUMBER == exp\_pkts - 1}).\\
\texttt{bresp\_ready} & 1 & B response path $\rightarrow$ Write ROB & Accepts the ordered BRESP and permits write-slot release.\\
\texttt{rd\_dbuf alloc/read/free} & TBD & Read ROB $\leftrightarrow$ Read Data Buffer & Handshake and release timing are \textbf{TBD}; a contiguous \texttt{exp\_pkts}-line range is allocated at ROB init and released as one unit.\\
\bottomrule
\end{longtable}

\section{Arbitration and backpressure}

The request validators and Burst Splitter retain the existing request
arbitration policy.  The selected ROB Free List and selected AXID Slot FIFO
provide transaction-allocation backpressure.  A full Slot FIFO affects only
its AXID; it does not imply global Free List exhaustion.  Conversely, an empty
Free List means all 256 transaction entries are allocated.

Read packet issuance is additionally controlled by pooled Read Data Buffer
availability and its independent 128-entry watermark.  Metadata RAM capacity
does not control packet issue.

\section{Reset and initialization sequence}
\label{sec:resetseq}

Reset is active-low, asserted asynchronously, and deasserted synchronously to
the CIF clock domain.  While asserted, it forces every structure below to its
initialized state; normal operation begins only after synchronous
deassertion.

On reset, each ROB instance independently initializes its 256-bit Free List
to free, clears Metadata RAM state, resets all 64 Slot FIFOs, and clears its
Retiring Register.  The write \texttt{wr\_line\_base} and \texttt{wr\_line\_count}
footprint fields are part of Metadata RAM state and are cleared with it; the
read-admission same-line hazard gate is combinational and holds no reset
state (\S\ref{sec:samelinehaz}).  The Burst Splitter clears its per-packet
address-layout state, \texttt{exp\_pkts}, \texttt{pkt\_idx}, and retained
\texttt{ROB\_INDEX}.
The CIF--MC async FIFOs (\S\ref{sec:cifmc}) reset to empty; each credit
counter resets to its FIFO depth.  No transaction crosses either seam until
both clock domains are out of reset.  The address-map / hash CSRs are
CIF-side and are loaded and locked before \texttt{init\_done}.

A soft reset first quiesces request acceptance and waits for packet FIFOs,
Burst Splitter activity, both ROB Slot-FIFO sets, both Retiring Registers, and
all externally visible AXI responses to drain.  A read held at the AR Metadata
FIFO head by the same-line hazard check (\S\ref{sec:samelinehaz}) is upstream
of ROB allocation; quiescing request acceptance stops feeding it and its
blocking writes drain through normal retirement, so no additional quiescence
condition is required.  The Read Data Buffer has a
separate drain/initialization contract; its exact interface and timing are
\textbf{TBD}.  Reset must not release a read \texttt{ROB\_INDEX} until the
final AXI read response and its Data Buffer release have completed.

\section{Error propagation}

Error packets are split and tagged in the same manner as ordinary
requests.  Every completion sets its packet-number bit in the selected ROB.
The write ROB aggregates packet status before its one BRESP; the read ROB
serializes status with its data.  The \texttt{status} encoding
(\texttt{0}/\texttt{1}/\texttt{2} $=$ OK/SLVERR/DECERR, codes \texttt{3..15}
reserved $\rightarrow$ DECERR) and the write worst-status precedence
(\texttt{DECERR} $>$ \texttt{SLVERR} $>$ \texttt{OKAY}) are frozen in
\S\ref{sec:cifmc}.
